\chapter{Závěr}
Cílem této bakalářské práce bylo porovnat vybrané databázové systémy z hlediska jejich vhodnosti pro ukládání a zpracování časových řad, a to jak z teoretického, tak z praktického hlediska, a na základě tohoto porovnání formulovat doporučení pro výběr vhodného systému.

V teoretické části (viz \ref{theoretical}) práce byla nejprve charakterizována podstata časových řad a přístupy k jejich ukládání, dále byly představeny základní typy databázových systémů (relační, dokumentové, sloupcové, párů klíč-hodnota a grafové) společně s jejich vnitřními datovými strukturami, zejména B-Tree a LSM stromem. 
Následně byly popsány dostupné nástroje a technologie použité pro přípravu testovacího prostředí: virtualizační a kontejnerizační platformy, orchestrační nástroje pro automatizaci konfigurace a monitorovací systémy.

V praktické části (viz \ref{practical}) bylo vybráno osm databázových systémů: Apache IoTDB, ClickHouse, \mbox{InfluxDB 3}, MariaDB, MariaDB ColumnStore, MongoDB, QuestDB a TimescaleDB. Tyto databázové systémy reprezentují různé přístupy k ukládání dat. 
Pro jejich porovnání bylo pomocí nástroje Ansible vytvořeno automatizované a reprodukovatelné testovací prostředí, ve kterém byly systémy nasazeny za srovnatelných podmínek na reálném datasetu poskytnutém Datovým centrem Ústeckého kraje. 
Systémy byly otestovány na třech hardwarových profilech sadou testů zaměřených na vkládání, bodové čtení, filtrování podle podmínky, agregaci a kontinuální (živý) zápis dat, přičemž průběh testů byl souběžně monitorován z hlediska využití CPU, paměti a diskových operací.

Z výsledků vyplynulo (viz \ref{vysledky}), že nejlépe si ve výkonnostních testech vedl ClickHouse, který dosáhl nejlepších nebo téměř nejlepších výsledků napříč všemi profily i typy dotazů, a to při současně nejjednodušší instalaci ze všech testovaných systémů, permisivní licenci Apache 2.0 a nativní podpoře horizontálního škálování. 
Jako nejsilnější alternativu lze označit TimescaleDB, který na výkonnějším hardwarovém profilu při kontinuálním zápisu ClickHouse dokonce mírně předčil a jako jediný ze všech systémů vykazoval u čtecích operací minimální citlivost na navýšení hardwarových prostředků. 
Naopak nejhorších výsledků v klasických testech dosahovalo MongoDB a nejnáročnější na instalaci a konfiguraci byla Apache IoTDB.

Na základě provedeného srovnání lze pro zpracování časových řad v prostředí Datového centra Ústeckého kraje jako celkově nejvhodnější doporučit databázový systém ClickHouse. 
Pokud by však konkrétní nasazení kladlo důraz spíše na časté zápisy menších objemů dat, například z jednotlivých IoT zařízení, představuje vhodnou alternativu TimescaleDB.

Je třeba zdůraznit, že uvedené závěry vycházejí z konkrétního datasetu telemetrie vozidel správy a údržby silnic a z testovacích scénářů navržených pro tento konkrétní případ užití, a proto je nelze bez dalšího zobecňovat na jakékoliv nasazení. 
Práce se rovněž omezila na open-source edice testovaných systémů, jednouzlové nasazení bez replikace či clusteringu a na testovací prostředí realizované formou LXD kontejnerů. 
Případné navazující práce by se mohly zaměřit na ověření chování systémů v distribuovaném nasazení, na porovnání komerčních Enterprise edic nebo na otestování dalších typů datových sad s odlišnou charakteristikou zápisu a čtení.

Cíle stanovené v úvodu práce byly naplněny. Byla vytvořena funkční automatizovaná testovací infrastruktura, byly provedeny a vyhodnoceny výkonnostní testy osmi databázových systémů a na základě získaných výsledků bylo formulováno jak konkrétní doporučení, tak obecnější kritéria pro výběr databázového systému vhodného pro zpracování časových řad.